\chapter{Logit and Probit}
\label{cap:logit}
% ── index ───────────────────────────────────────────────────────
\index{logit|textbf}
\index{probit|textbf}
\index{logit@\texttt{logit()}|textbf}
\index{probit@\texttt{probit()}|textbf}
\index{maximum likelihood|textbf}
\index{marginal effects|textbf}
\index{margins@\texttt{margins()}|textbf}
\index{odds ratio}

% ─────────────────────────────────────────────────────────────────────────────
\section{The latent variable model}

When the dependent variable is binary ($y \in \{0,1\}$), the linear model
produces predictions outside the $[0,1]$ interval and heteroscedastic errors by
construction. The solution is to model the probability of $y=1$ via a
cumulative distribution function:

\[
  P(y_i = 1 \mid x_i) = F(x_i'\beta)
\]

The motivation comes from the latent variable model:
\[
  y_i^* = x_i'\beta + \varepsilon_i, \qquad y_i = \mathbf{1}[y_i^* > 0]
\]

The distribution of $\varepsilon$ determines the model:

\begin{center}
\begin{tabular}{lll}
\toprule
\textbf{Model} & \textbf{Distribution of $\varepsilon$} & \textbf{$F(\cdot)$} \\
\midrule
Logit  & Logistic$(0, \pi^2/3)$ & $\Lambda(z) = \dfrac{1}{1 + e^{-z}}$ \\[8pt]
Probit & Normal$(0, 1)$          & $\Phi(z) = \displaystyle\int_{-\infty}^{z} \phi(t)\,dt$ \\
\bottomrule
\end{tabular}
\end{center}

Estimation is by Maximum Likelihood (MLE). The log-likelihood is:
\[
  \ell(\beta) = \sum_{i=1}^{n}
    \left[y_i \ln F(x_i'\beta) + (1-y_i)\ln\bigl(1 - F(x_i'\beta)\bigr)\right]
\]

% ─────────────────────────────────────────────────────────────────────────────
\section{Interpretation of coefficients}

Logit and Probit coefficients \textbf{are not marginal effects} --- they are
parameters of the latent variable, without a directly interpretable unit.
For continuous $x_j$, the marginal effect on the probability is:

\[
  \frac{\partial P(y=1 \mid x)}{\partial x_j}
  = f(x'\beta)\,\beta_j
\]

where $f = F'$ is the density of the chosen distribution. Since $f(x'\beta)$
varies across observations, the \textbf{Average Marginal Effect} (AME) is reported:

\[
  \mathrm{AME}_j = \frac{1}{n}\sum_{i=1}^{n} f(x_i'\beta)\,\beta_j
\]

\begin{nota}
  The relationship between Logit and Probit coefficients is approximately:
  \[
    \hat\beta_{\mathrm{Logit}} \approx \frac{\pi}{\sqrt{3}}\,\hat\beta_{\mathrm{Probit}}
    \approx 1{,}81\,\hat\beta_{\mathrm{Probit}}
  \]
  The AMEs of the two models are typically very close, even with
  coefficients on different scales.
\end{nota}

% ─────────────────────────────────────────────────────────────────────────────
\section{DGP Verification}

\subsection*{MLE consistency}

Unlike the linear models in previous chapters, Logit and Probit
\emph{do not admit exact recovery} of parameters. The binary variable $y$
is generated by Bernoulli sampling --- there is irreducible randomness even
without adding extra noise. What is verified here is \textbf{consistency}:
with $n=1000$ the MLE should converge close to the true parameters.

The DGP uses the logistic parametrization with $\beta_0 = -1{,}0$ and $\beta_1 = 2{,}0$:
\[
  x_i \sim \mathcal{N}(0,1), \qquad
  p_i = \Lambda(-1{,}0 + 2{,}0\,x_i), \qquad
  y_i \sim \mathrm{Bernoulli}(p_i)
\]

Individual Bernoulli is generated by $y_i = \mathbf{1}[u_i < p_i]$, with
$u_i \sim \mathrm{Uniform}(0,1)$:

\begin{lstlisting}[caption={Logistic DGP --- Logit and Probit simultaneously}]
seed(42)
input df
id
1
end
for i in 1..=10 { df = append(df, df) }
df |> filter(_n <= 1000)
generate df x  = rnormal()
generate df lp = -1.0 + 2.0*x
generate df p  = 1 / (1 + exp(-lp))
generate df u  = runiform()
generate df y  = u < p
let m_l = logit(y ~ x, df)
let m_p = probit(y ~ x, df)
esttab(m_l, m_p)
\end{lstlisting}

\begin{lstlisting}[language={}, basicstyle=\ttfamily\small,
  frame=none, numbers=none, backgroundcolor=\color{codbg},
  xleftmargin=0pt, xrightmargin=0pt, breaklines=false]
                            (1)              (2)
                          Logit           Probit
────────────────────────────────────────────────
const                -0.6772***       -0.4222***
                       (0.1326)         (0.0821)
x                     1.3329***        0.8324***
                       (0.2308)         (0.1427)
────────────────────────────────────────────────
N                          1000             1000
────────────────────────────────────────────────
* p<0.10  ** p<0.05  *** p<0.01
\end{lstlisting}

The true parameters are $(\beta_0, \beta_1) = (-1{,}0;\; 2{,}0)$ for
the Logit. Since \lstinline{rnormal()} generates $U(0,1)$ instead of $\mathcal{N}(0,1)$
(see ch.~\ref{cap:ols}), $x \in [0,1]$ and the linear predictor $lp \in [-1,1]$
varies little --- which reduces Fisher information and attenuates estimates
relative to the true DGP. Probit estimates on the normal scale; the values
$(-0{,}42;\; 0{,}83)$ are within the 95\% CI of the probit equivalents
$(-0{,}55;\; 1{,}10)$.

\subsection*{Average marginal effects}

Coefficients are only parameters of the latent variable. For
economic interpretation, use \lstinline{margins}:

\begin{lstlisting}[caption={Logit AME}]
margins(m_l, df)
\end{lstlisting}

\begin{lstlisting}[language={}, basicstyle=\ttfamily\small,
  frame=none, numbers=none, backgroundcolor=\color{codbg},
  xleftmargin=0pt, xrightmargin=0pt, breaklines=false]
════════════════════════════════════════════════════════════
 Average Marginal Effects — LOGIT
════════════════════════════════════════════════════════════
Variable                dy/dx   Std.Err.        z    P>|z|
────────────────────────────────────────────────────────────
x                    0.321859   0.035576    9.047   0.0000 ***
────────────────────────────────────────────────────────────
n = 1000
════════════════════════════════════════════════════════════
\end{lstlisting}

A unit increase in $x$ raises the probability of $y=1$ by
approximately $0{,}32$ percentage points, on average over the sample.

% ─────────────────────────────────────────────────────────────────────────────
\section{Syntax}

\begin{lstlisting}
logit(y ~ x1 + x2, df)
probit(y ~ x1 + x2, df)
margins(model, df)
\end{lstlisting}

\begin{lstlisting}[caption={Typical workflow --- discrete choice}]
load "data.csv" as df

let m_l = logit(employed ~ educ + exp + sex, df)
let m_p = probit(employed ~ educ + exp + sex, df)

esttab(m_l, m_p)
margins(m_l, df)
\end{lstlisting}

\section{Predicted probabilities}

\begin{lstlisting}
let m = logit(y ~ x, df)
let phat = predict(m, "pr")     // P(y=1|x) for each obs
let yhat = predict(m, "class")  // predicted class (0 or 1)
\end{lstlisting}

\begin{nota}
  Use \lstinline{predict(m, "class")} with default threshold of 0{,}5.
  In imbalanced data (low prevalence), adjust the threshold with
  \lstinline{predict(m, "class", threshold=0.3)}.
\end{nota}
